\chapter{Resultados y análisis}\label{chap:resultados}
Este capítulo presenta la evidencia empírica del sistema propuesto a partir del modelo final entrenado con 84 descriptores y selección automática de 20 variables mediante \texttt{SelectKBest}. Se reportan métricas de discriminación, calibración y robustez, así como el desempeño del detector de anomalías. Todas las estimaciones se obtienen bajo el régimen de validación anidada descrito en el Capítulo~\ref{chap:metodologia}, con control explícito de semillas y artefactos versionados para asegurar reproducibilidad \cite{Tsamardinos2022,Vabalas2019,Pedregosa2011}.

\section{Desempeño del modelo y separabilidad de las características}
El ensamble final—combinación convexa de regresión logística penalizada, SVM con núcleo RBF, bosque aleatorio y potenciación de gradiente—alcanza una exactitud en prueba de \(0{,}951\) y un ROC-AUC de \(0{,}989\). La Tabla~\ref{tab:metricas-globales} resume las métricas globales, incluyendo la media y desviación estándar de exactitud en validación cruzada estratificada de 5 pliegues.

\begin{table}[H]
    \centering
    \caption{Métricas del modelo final en el conjunto de prueba y en validación cruzada.}
    \label{tab:metricas-globales}
    \begin{tabular}{lccc}
        \toprule
        Modelo & Exactitud test & AUC ROC & CV Accuracy (media $\pm$ DE) \\
        \midrule
        Ensamble final & 0,951 & 0,989 & $0,934 \pm 0,020$ \\
        \bottomrule
    \end{tabular}
\end{table}

\begin{figure}[H]
    \centering
    \includegraphics[width=0.95\textwidth]{figuras/feature_distributions.png}
    \caption{Distribuciones empíricas de las características más discriminativas. Cada panel compara clases sanas (azul) y falciformes (naranja) en pares de atributos; las líneas punteadas marcan medias por clase.}
    \label{fig:feature-distributions}
\end{figure}

La Figura~\ref{fig:feature-distributions} evidencia que varias combinaciones de atributos producen fronteras casi lineales entre cohortes, coherentes con el buen desempeño de la regresión logística dentro del ensamble. En particular, los perfiles derivados de coocurrencias (p.\,ej., correlación GLCM), la energía en anillos de Fourier y la energía de detalles wavelet separan con baja solapación las distribuciones de ambas clases, lo que reduce la varianza del estimador y refuerza la estabilidad del sistema ante pequeñas perturbaciones \cite{Barburiceanu2021,Amann2019,Alzubaidi2020}. Este comportamiento avala el diseño del vector de características: descriptores texturales y frecuenciales con base física (formación de franjas y contenido espectral) que, al combinarse, capturan anisotropías y granularidad propias de los hologramas de eritrocitos \cite{OConnor2020}.

\section{Matriz de confusión y curva ROC}
\begin{figure}[H]
    \centering
    \includegraphics[width=0.75\textwidth]{figuras/confusion_matrix.png}
    \caption{Matriz de confusión normalizada para el conjunto de prueba.}
    \label{fig:confusion}
\end{figure}

La matriz de confusión de la Figura~\ref{fig:confusion} muestra un equilibrio clínicamente favorable: sensibilidad del 1{,}00 para la clase falciforme (tasa de verdaderos positivos) y especificidad del 0{,}90 para la clase sana, consistente con una política de umbral que prioriza la detección de casos patológicos. Esta asimetría controlada responde al criterio costo–sensible adoptado en el capítulo metodológico, en el que se busca minimizar falsos negativos dada su mayor implicación clínica \cite{Tsamardinos2022,Vabalas2019}. La estabilidad de estas tasas se corrobora en la validación cruzada, con dispersión acotada de la precisión balanceada.

\begin{figure}[H]
    \centering
    \includegraphics[width=0.75\textwidth]{figuras/roc_curve.png}
    \caption{Curva ROC del modelo final con el área bajo la curva correspondiente.}
    \label{fig:roc}
\end{figure}

La curva ROC de la Figura~\ref{fig:roc} confirma una discriminación alta a lo largo de todo el rango de umbrales, con \(\mathrm{AUC}=0{,}989\). La proximidad al vértice superior izquierdo indica que, para umbrales intermedios, la sensibilidad y la especificidad superan simultáneamente el 95\,\%. Este resultado es coherente con la geometría observada en las distribuciones de la Figura~\ref{fig:feature-distributions}, donde el solapamiento interclase es reducido, y con la combinación de modelos del ensamble, que amortigua sesgos particulares de cada base \cite{Pedregosa2011,Friedman2001}.

\section{Detección de anomalías y estructura del espacio de características}
\begin{table}[H]
    \centering
    \caption{Resultados de la detección de anomalías en el repositorio externo.}
    \label{tab:anomalias}
    \begin{tabular}{lccc}
        \toprule
        Categoría externa & Imágenes analizadas & Distancia media & Tasa de anomalías \\
        \midrule
        \textit{8\_Poaceae} & 10 & 85,3 & 100~\% \\
        \textit{2\_Aspergillus} & 10 & 142,1 & 100~\% \\
        \textit{3\_Diatom} & 10 & 105,2 & 100~\% \\
        \bottomrule
    \end{tabular}
\end{table}

La Tabla~\ref{tab:anomalias} resume el desempeño del detector de Mahalanobis, ajustado sobre el mismo subespacio de características del clasificador y con umbral operativo fijado en el percentil 95 de las distancias del entrenamiento (\(9{,}89\)). Las categorías externas—organismos ajenos a la hematología—presentan distancias medias dos órdenes de magnitud por encima del umbral, con una tasa de detección del 100\,\%. Este margen amplio respalda la utilidad del detector como salvaguarda: ante entradas fuera de dominio, la alerta es inequívoca, lo que favorece la abstención informada y la derivación a revisión humana \cite{Mahalanobis1936,Rousseeuw2019,Montero2024}.

\begin{figure}[H]
    \centering
    \includegraphics[width=0.8\textwidth]{figuras/tsne_mahalanobis.png}
    \caption{Proyección t-SNE del espacio de características, coloreada por distancia de Mahalanobis.}
    \label{fig:tsne-mahalanobis}
\end{figure}

La proyección de la Figura~\ref{fig:tsne-mahalanobis} ilustra la estructura geométrica del espacio de características: los clústeres correspondientes a sanos y falciformes forman dos conglomerados dominantes con gradientes de distancia de Mahalanobis crecientes hacia regiones periféricas. La concentración de símbolos con coloraciones intensas fuera de los núcleos de clase indica que el detector es sensible a desviaciones multivariadas coherentes con la covarianza estimada, y que no responde a variaciones marginales aisladas. Esta visualización no participa del entrenamiento, pero corrobora la separación inducida por los descriptores y justifica el umbral adoptado.

\section{Incertidumbre, calibración y robustez estadística}
\begin{figure}[H]
    \centering
    \includegraphics[width=0.95\textwidth]{figuras/performance_uncertainty.png}
    \caption{Exactitud con intervalos bootstrap y curva de confiabilidad del modelo final.}
    \label{fig:performance-uncertainty}
\end{figure}

La Figura~\ref{fig:performance-uncertainty} sintetiza dos aspectos clave. A la izquierda, el intervalo de confianza al 95\,\% para la exactitud estimado por \emph{bootstrap} estratificado (2000 remuestreos) muestra una dispersión acotada a pesar del tamaño muestral limitado, en concordancia con las recomendaciones para evaluación en pequeños \(n\) \cite{Vabalas2019,Tsamardinos2022}. A la derecha, la curva de confiabilidad permanece próxima a la diagonal ideal con una pérdida de Brier de \(0{,}039\); esto indica que las probabilidades emitidas por el ensamble son informativas y coherentes con las frecuencias observadas, requisito cuando las salidas se emplean para priorización clínica \cite{Pedregosa2011}. La calibración se obtuvo en un pliegue independiente del utilizado para búsqueda de hiperparámetros, preservando la independencia entre etapas y evitando optimismo en la estimación \cite{Tsamardinos2022}.

\section{Análisis de errores y discusión}
El conjunto de prueba registró catorce clasificaciones incorrectas, repartidas entre falsos positivos y negativos. Los falsos negativos se concentran en hologramas de células falciformes con contraste reducido y contenido espectral de altas frecuencias cercano al de eritrocitos sanos; los falsos positivos aparecen en hologramas sanos con patrones de interferencia elongados. En ambos grupos, las contribuciones de características sitúan a la solidez y a la energía wavelet cerca de los umbrales de decisión, lo que sugiere incrementar, en futuras iteraciones, la representación de estas zonas grises dentro del entrenamiento y explorar ajustes finos de los parámetros de Gabor y de la cuantización GLCM para mejorar sensibilidad local \cite{Barburiceanu2021,Alzubaidi2020}. 

En conjunto, los resultados corroboran la hipótesis de trabajo: un conjunto de descriptores físicos interpretables, combinado con un ensamble de clasificadores clásicos, alcanza métricas elevadas sin recurrir a reconstrucción del campo complejo. La discriminación (\(\mathrm{AUC}=0{,}989\)), la calibración (Brier \(=0{,}039\)) y la detección inequívoca de entradas fuera de dominio sustentan la viabilidad del enfoque para escenarios de cribado, donde la trazabilidad de decisiones y la comunicación de incertidumbre son indispensables \cite{Tsamardinos2022,Montero2024}. La evidencia sugiere líneas de mejora incremental—enriquecimiento de casos límite y expansión multicéntrica del dataset—compatibles con la arquitectura y el régimen de validación actuales.
